\subsubsection{Bilevel Continual Learning (BCL-Dual)}

\gls{bcl} \cite{ll_bcl_2020} is a bi-level meta-learning approach to \gls{cl} that
frames continual learning as a nested optimisation problem. The BCL-Dual variant
maintains two memory streams and alternates between inner and outer optimisation loops to
learn parameters that generalise across tasks. \gls{bcl} requires task identity
during training to manage per-task memory partitions, but does not require task
identity at inference time, making it compatible with both \gls{til} and
\gls{cil} evaluation settings.

\paragraph{Bi-level Optimisation}

The motivating objective for \gls{bcl} is the nested optimisation:
%
\begin{equation}\label{eqn:bcl_bilevel}
    \min_{\theta} \quad \mathcal{L}^{\text{meta}}(\theta, \phi^*(\theta))
    \quad \text{s.t.} \quad
    \phi^*(\theta) = \operatorname*{arg\,min}_{\phi} \;
    \mathcal{L}^{\text{train}}(\phi, \theta)
\end{equation}
%
where $\theta$ are the persistent main-model parameters and $\phi$ are fast
weights initialised from $\theta$ and adapted on the current task. Solving this
exactly requires differentiating through the inner optimisation — i.e.\
computing second-order gradients of $\mathcal{L}^{\text{meta}}$ with respect
to $\theta$ via $\phi^*(\theta)$. \gls{bcl} adopts a first-order approximation
analogous to \gls{fomaml}: fast weights $\phi$ are initialised from $\theta$,
updated by two sequential gradient steps, and then discarded; $\theta$ is
updated using the resulting parameter displacement. $K$ such iterations are
performed per observation (typically $K=5$).

\subparagraph{Inner Step.} At the start of each iteration, a snapshot
$\theta$ of the current main-model parameters is taken and fast weights are
initialised as $\phi \leftarrow \theta$. The network then receives both the
current task minibatch and a replay sample drawn from the working memory buffer,
and $\phi$ is updated via a gradient step on $\mathcal{L}^{\text{train}}$:
%
\begin{equation}\label{eqn:bcl_inner}
    \mathcal{L}^{\text{train}}(\phi) =
    \lambda_{\text{cls}} \, \mathcal{L}_{\text{CE}}(x, y)
    + \mathcal{L}_{\text{CE}}(x^{\text{mem}}, y^{\text{mem}})
    + \lambda_{\text{mem}} \,
      D_{\text{KL}}\!\left(
          \sigma\!\left(\tfrac{f_\phi(x^{\text{mem}})}{\tau}\right)
          \;\Big\|\;
          \hat{p}^{\text{mem}}
      \right)
\end{equation}
%
where $x^{\text{mem}}, y^{\text{mem}}$ are the replayed inputs and labels,
$\hat{p}^{\text{mem}}$ are the soft-target predictions cached in the working
memory at the time each sample was last observed, $\tau$ is the distillation
temperature, $\lambda_{\text{cls}}$ weights the current-task classification
term, and $\lambda_{\text{mem}}$ weights the distillation regulariser:
%
\begin{equation}
    \phi \leftarrow \phi - \alpha \nabla_\phi \mathcal{L}^{\text{train}}(\phi)
\end{equation}

\subparagraph{Outer Step.} A second forward pass on the meta memory buffer
evaluates $\mathcal{L}^{\text{meta}}$ — a cross-entropy loss on held-out samples
from all tasks seen so far — at the current fast-weight state $\phi$. A further
gradient step updates $\phi$ to $\phi'$:
%
\begin{equation}
    \phi' = \phi - \alpha \nabla_{\phi} \mathcal{L}^{\text{meta}}(\phi)
\end{equation}
%
The main-model parameters are then updated using the total displacement from
$\theta$ to $\phi'$, scaled by $\beta$, and the fast weights are discarded:
%
\begin{equation}\label{eqn:bcl_outer}
    \theta \leftarrow \theta + \beta \left( \phi' - \theta \right)
\end{equation}
%
The scalar $\beta$ controls how far $\theta$ moves towards $\phi'$. When
$\beta < 1$ the update is conservative; when $\beta = 1$ $\theta$ is set equal
to $\phi'$; when $\beta > 1$ it extrapolates beyond $\phi'$. In practice
$\beta$ is a hyperparameter: larger values improve plasticity in \gls{til}
settings, whilst lower values are preferred for \gls{cil} where more
conservative updates reduce interference between classes. Because
$\mathcal{L}^{\text{meta}}$ is evaluated at the fast-adapted $\phi$ rather
than at $\theta$, the outer gradient implicitly accounts for the inner
adaptation without requiring second-order derivatives.

\paragraph{Dual Memory Streams}

The BCL-Dual variant maintains two memory buffers with a fixed total capacity
$M$ that is partitioned evenly across tasks. Both operate as ring buffers once
full, consistent with the online \gls{cl} setting. By default, 80\% of each
task's allocation is reserved for the working memory and 20\% for the meta
memory.

\subparagraph{Working Memory (Inner Loop).} The working memory stores, for each
observed sample, the raw input $x$, the ground-truth label $y$, and the
soft-target prediction $\hat{p} = \sigma(f_\theta(x)/\tau)$ produced by the
network at the time the sample was written. The soft targets are refreshed
whenever the buffer slot is overwritten, so they reflect reasonably recent
network state. This buffer is sampled during the inner optimisation loop to
supply the replay cross-entropy term and the KL distillation term in
Equation~\ref{eqn:bcl_inner}. Because stored soft targets encode the
network's belief at write time, the distillation loss penalises large drifts in
predictive distribution on past inputs, acting as an implicit regulariser
against catastrophic forgetting.

\subparagraph{Meta Memory (Outer Loop).} The meta memory stores only the raw
input $x$ and the ground-truth label $y$ — no soft targets are cached. It acts
as a held-out validation set for the outer loop, providing the cross-entropy
signal used to compute the meta gradient in
Equation~\ref{eqn:bcl_outer}. Because this buffer is never used in the inner
optimisation, its labels remain clean supervision targets unseen by the inner
loop, allowing the outer gradient to measure generalisation across tasks rather
than fit to the current training stream.

The meta-objective $\mathcal{L}^{\text{meta}}$ evaluated on this buffer
encourages the network to reach a parameter configuration from which fast
adaptation to new tasks is possible without forgetting prior ones, in the
spirit of MAML-style meta-learning.

\paragraph{Relationship to CTN}

Both \gls{bcl} and \gls{ctn} employ dual memory streams and bi-level parameter
updates; the key distinction is in what each level optimises. In \gls{ctn}, the
two levels correspond to backbone and controller parameters, with the controller
learning task-specific \gls{film} transformations. In \gls{bcl}, the two levels
correspond to task-specific base parameters and task-agnostic meta parameters,
with the outer loop explicitly optimising for cross-task generalisation. This
makes \gls{bcl} more closely related to meta-learning approaches such as MAML,
whilst \gls{ctn} is more closely related to modular and feature-transformation
methods.